\documentclass[10pt,a4paper]{amsart}
\usepackage[margin=19mm]{geometry}
\usepackage{amsmath,amssymb,mathtools}
\usepackage[scr=rsfs,cal=euler]{mathalfa}
\usepackage{xfrac}
\usepackage[unicode,hidelinks,bookmarks=false]{hyperref}
\newcommand{\ie}{i.e.\ }
\newcommand{\cf}{cf.\ }
\newcommand{\resp}{respectively\ }
\newcommand{\Subsection}{\subsection}
\providecommand{\nobreakdash}{\nobreak\hbox{-}}
\setlength{\emergencystretch}{2em}
\multlinegap0pt
\usepackage{amssymb}
\usepackage[shortlabels]{enumitem}
\usepackage{float}
\newtheorem{theorem}{Theorem}
\title{Analytic patch trees: branch interface inheritance and fractal dimension fields}
\author{Henk Mulder}
\date{Source: arXiv:2606.06400v1 (CC BY 4.0)}
\begin{document}
\maketitle

\noindent\emph{Source and licence.} Analytic patch trees: branch interface inheritance and fractal dimension fields. Version arXiv:2606.06400v1 (CC BY 4.0), distributed by arXiv under the Creative Commons Attribution 4.0 International licence, http://creativecommons.org/licenses/by/4.0/, as recorded in the arXiv licence metadata for this exact version and shown on the abstract page https://arxiv.org/abs/2606.06400v1. Full licence notice: \texttt{licenses/arxiv-2606.06400-CC-BY-4.0.txt}.\par\smallskip

\noindent\textit{Editorial scope.} Counted unit: the article's theorem cor:canonical\_selfsimilar (source0.tex lines 694--741). The article's complete original proof of it is delivered with it (source0.tex lines 742--828, one \texttt{proof} environment of 87 lines). No mathematics is written, repaired or re-derived: the statement and the proof are the article's own lines, in the article's own notation, and no retained line is substituted or re-flowed.  No further source-local context is needed: every cross-reference of every retained line resolves inside the two delivered spans.  Nothing was omitted from any retained span, so the package carries no omission record.  No retained line contains a citation, so the wrapper carries no bibliography and no bibliography entry.  No retained line contains a figure environment, an image-inclusion command or drawing code, so no image is delivered and the package declares no asset dependency.  Editorial formatting only, in addition: an AMS \texttt{amsart} preamble that restates the article's own declarations and macros used by the retained lines, namely the environments \texttt{theorem} on separate counters, together with the standard packages those lines need.  The retained environments print in this document as Theorem 1 in order of appearance, whereas the article prints the same items with the numbering of its own full document.  The complete native source of the article as distributed in the arXiv e-print for this exact version is bundled unchanged as \texttt{sources/202185/source0.tex}.
\begin{theorem}[Canonical Self-Similar Conformal Tree]
\label{cor:canonical_selfsimilar}

Let the conformal generator field be

\[
    \phi(z)=A+kz,
\]

with $A,k\in\mathbb{C}$ constant. Then the corresponding realised
patch map satisfies

\[
    F(z+iL)
    =
    aF(z)+b,
\]

where

\[
    a=e^{ikL},
    \qquad
    b=(1-a)F_0.
\]

Consequently, the interface evolution operator acts by a similarity
transformation on the inherited interface family. The resulting
conformal patch tree is therefore geometrically self-similar, with
successive interfaces belonging to a single similarity class.

In particular, writing

\[
k=\alpha+i\beta,
\]

the similarity is contractive whenever

\[
\beta>0.
\]

\[
    F_*=\frac{b}{1-a}=F_0.
\]

\end{theorem}

\begin{proof}
Since $\phi(z)=A+kz$, we have

\[
    F'(z)=e^{\phi(z)}=e^A e^{kz}=Ce^{kz}.
\]

Integrating gives

\[
    F(z)=De^{kz}+F_0,
\]

for constants $D,F_0\in\mathbb{C}$. Therefore

\[
    F(z+iL)-F_0
    =
    De^{k(z+iL)}
    =
    e^{ikL}De^{kz}
    =
    e^{ikL}\bigl(F(z)-F_0\bigr).
\]

Hence

\[
    F(z+iL)
    =
    e^{ikL}F(z)+(1-e^{ikL})F_0.
\]

Setting

\[
    a=e^{ikL},
    \qquad
    b=(1-a)F_0,
\]

we obtain

\[
    F(z+iL)=aF(z)+b.
\]

This is a similarity transformation of the realised plane whenever
$a\neq 0$, which holds because $a=e^{ikL}$. Thus the tip interface is
similar to the base interface, and the interface evolution operator
preserves the similarity class of the inherited interfaces. Repeating
the same argument generation by generation shows that all successive
interfaces are related by iterates of the same similarity map.

If $|a|<1$, these iterates converge to the fixed point $F_*$ satisfying

\[
    F_*=aF_*+b.
\]

Solving gives

\[
    F_*=\frac{b}{1-a}=F_0.
\]

Since

\[
    |a|=|e^{ikL}|,
\]

the contraction condition is determined by the imaginary part of $k$.
For $k=\alpha+i\beta$,

\[
    ikL=i\alpha L-\beta L,
\]

and therefore

\[
    |a|=e^{-\beta L}.
\]

Thus the similarity is contractive precisely when $\beta>0$.
\end{proof}

\end{document}
